% Final page of the QQ / GF(2) orbit atlas.
% Input after tree_atlas_style.tex; graphicx and arrows.meta are also required.
% Asset paths are relative to this directory, as in the enclosing atlas.
% Vertex correspondence is exact modulo common lineality:
% A = QQ 34 = GF2 1, B = QQ 9 = GF2 2,
% C = QQ 40 = GF2 13, D = QQ 6 = GF2 14; O = GF2 3.
% The QQ orbit labels below were checked against 02_coarsest_cells.csv.
\begingroup
\definecolor{ftQQBlue}{HTML}{397997}
\definecolor{ftQQTeal}{HTML}{268375}
\definecolor{ftQQPurple}{HTML}{81629C}
\definecolor{ftGFRed}{HTML}{B55F61}
\definecolor{ftGFTeal}{HTML}{398A77}
\definecolor{ftGFPurple}{HTML}{81629C}
\definecolor{ftGFGold}{HTML}{BA8A38}
\tikzset{
  ft vertex/.style={circle,draw=ink!55,fill=white,line width=.9pt,
    minimum size=6.2mm,inner sep=0pt,font=\fontsize{14}{14}\selectfont\bfseries,text=ink},
  ft central/.style={ft vertex,draw=heartred,fill=heartred!9,line width=1.3pt},
  ft edge/.style={draw=ink!65,line width=1.35pt,line cap=round},
  ft hidden/.style={ft edge,densely dashed,draw=ink!45},
  ft centre edge/.style={draw=heartred!55,line width=1.25pt,line cap=round},
  ft face label/.style={fill=white,fill opacity=.95,text opacity=1,rounded corners=1.5pt,
    inner sep=1.0mm,font=\fontsize{13}{14}\selectfont\bfseries},
  ft small/.style={text=muted,font=\fontsize{11}{13}\selectfont,align=center},
  ft tree caption/.style={text=muted,font=\fontsize{10.5}{12}\selectfont,align=center},
  ft title/.style={text=ink,font=\fontsize{15.5}{17}\selectfont\bfseries,align=center},
  ft arrow/.style={line width=1.65pt,line cap=round,line join=round,
    -{Stealth[length=3.2mm,width=2.4mm]}},
  ft leader/.style={line width=1.25pt,line cap=round},
  ft underlay/.style={draw=white,line width=4.7pt,line cap=round,line join=round},
}
% All miniature triangles are labelled: transparency in the tetrahedron
% projection cannot make their incidences ambiguous.
\newcommand{\FTTriangle}[6]{%
  \begin{scope}[shift={(#1,#2)}]
    \path[fill=#3!18,draw=#3,line width=1.3pt,line join=round]
      (-.70,-.42)--(.70,-.42)--(0,.70)--cycle;
    \foreach \pos/\lbl in {(-.70,-.42)/#4,(.70,-.42)/#5,(0,.70)/#6}{%
      \node[circle,fill=white,draw=#3,line width=.65pt,minimum size=4.3mm,
        inner sep=0pt,font=\fontsize{10}{10}\selectfont\bfseries,text=ink]
        at \pos {$\lbl$};}
  \end{scope}}
\noindent\begin{tikzpicture}[x=1cm,y=1cm]
  \path[use as bounding box] (0,0) rectangle (39.5,24.6);
  \node[font=\fontsize{25}{28}\selectfont\bfseries,text=ink] at (19.75,24.0)
    {The field-dependent tetrahedra};
  \node[font=\fontsize{15}{17}\selectfont,text=muted] at (19.75,23.05)
    {In each of three Dressian tetrahedra, four boundary faces are replaced by six centre-to-edge triangles.};
  \node[font=\fontsize{18}{20}\selectfont\bfseries,text=ink] at (4.65,21.9)
    {$\mathbb Q$ only\;\textemdash\;3 maximal-cell orbits};
  \node[font=\fontsize{18}{20}\selectfont\bfseries,text=ink] at (34.8,21.9)
    {$\operatorname{GF}(2)$ only\;\textemdash\;4 maximal-cell orbits};

  % Left: the three QQ-only orbits.  The input trees use the census's
  % canonical orbit representatives, while the labelled triangles identify
  % the members of the displayed tetrahedron.
  \path[rounded corners=3pt,fill=ftQQBlue!4,draw=ftQQBlue!38,line width=.8pt]
    (0.05,15.55) rectangle (9.25,21.25);
  \node[ft title] at (4.65,20.68) {$\mathbb Q$ orbit 24\qquad size 6};
  \node at (2.65,18.58) {\resizebox{4.25cm}{!}{% Generated by export_field_atlas_assets.jl from the maintained census.
\begingroup
\definecolor{g2TaxonBlue}{HTML}{1676A3}
\definecolor{g2TaxonYellow}{HTML}{E39A24}
\definecolor{g2HeartRed}{HTML}{D7443E}
\definecolor{g2EdgeBlue}{HTML}{205B7A}
\tikzset{
  g2tree pendant/.style={draw=black!43,line width=.6pt,line cap=round},
  g2tree internal/.style={draw=g2EdgeBlue,line width=2pt,line cap=round},
  g2tree vertex/.style={circle,fill=g2EdgeBlue,draw=white,line width=.4pt,inner sep=0pt,minimum size=4.8pt},
  g2tree dark/.style={fill=g2TaxonBlue,draw=white,line width=.55pt,text=white,inner sep=0pt},
  g2tree light/.style={fill=g2TaxonYellow,draw=white,line width=.55pt,text=white,inner sep=0pt},
  g2tree weight/.style={circle,fill=white,draw=none,text=black,inner sep=.15pt,minimum size=10.0pt},
  g2tree aligned pendant/.style={fill=white,text=g2EdgeBlue,inner sep=.7pt,sloped,above=1pt,font={\sffamily\fontsize{9.25}{9.75}\selectfont}},
  g2tree aligned internal/.style={fill=white,text=g2EdgeBlue,inner sep=1pt,sloped,above=1pt,font={\sffamily\fontsize{9.25}{9.75}\selectfont\bfseries}},
}
\begin{tikzpicture}[x=1.0cm,y=1.0cm,baseline=(current bounding box.center)]
\coordinate (v1) at (1.65748,1.32180);
\coordinate (v2) at (-2.06685,-0.47174);
\coordinate (v3) at (0.91983,-1.91005);
\coordinate (v4) at (-0.91983,-1.91005);
\coordinate (v5) at (-1.65748,1.32180);
\coordinate (v6) at (2.06685,-0.47174);
\coordinate (v7) at (0.00000,2.12000);
\coordinate (v8) at (0.22011,-0.63222);
\coordinate (v9) at (-0.70723,0.21627);
\coordinate (v10) at (-0.03678,0.41753);
\coordinate (v11) at (0.52389,-0.00158);
\draw[g2tree internal] (v8)--(v11);
\draw[g2tree internal] (v9)--(v10);
\draw[g2tree internal] (v10)--(v11);
\draw[g2tree pendant] (v1)--(v11);
\draw[g2tree pendant] (v2)--(v9);
\draw[g2tree pendant] (v3)--(v8);
\draw[g2tree pendant] (v4)--(v8);
\draw[g2tree pendant] (v5)--(v9);
\draw[g2tree pendant] (v6)--(v11);
\draw[g2tree pendant] (v7)--(v10);
\node[g2tree weight,font={\sffamily\fontsize{9.25}{9.75}\selectfont}] at (0.37200,-0.31690) {2};
\node[g2tree weight,font={\sffamily\fontsize{9.25}{9.75}\selectfont}] at (-0.37200,0.31690) {1};
\node[g2tree weight,font={\sffamily\fontsize{9.25}{9.75}\selectfont}] at (0.24356,0.20798) {1};
\node[g2tree weight,font={\sffamily\fontsize{9.25}{9.75}\selectfont}] at (1.09069,0.66011) {0};
\node[g2tree weight,font={\sffamily\fontsize{9.25}{9.75}\selectfont}] at (-1.38704,-0.12774) {0};
\node[g2tree weight,font={\sffamily\fontsize{9.25}{9.75}\selectfont}] at (0.56997,-1.27114) {0};
\node[g2tree weight,font={\sffamily\fontsize{9.25}{9.75}\selectfont}] at (-0.34986,-1.27114) {0};
\node[g2tree weight,font={\sffamily\fontsize{9.25}{9.75}\selectfont}] at (-1.18235,0.76904) {0};
\node[g2tree weight,font={\sffamily\fontsize{9.25}{9.75}\selectfont}] at (1.29537,-0.23666) {0};
\node[g2tree weight,font={\sffamily\fontsize{9.25}{9.75}\selectfont}] at (-0.01839,1.26876) {1};
\node[g2tree vertex] at (v8) {};
\node[g2tree vertex] at (v9) {};
\node[g2tree vertex] at (v10) {};
\node[g2tree vertex] at (v11) {};
\node[g2tree dark,rectangle,minimum size=6.21mm,font={\sffamily\fontsize{12}{12}\selectfont\bfseries}] at (v1) {1};
\begin{scope}[shift={(v2)},scale=1.35,transform shape]
\path[fill=g2TaxonBlue,draw=white,line width=.55pt,line join=round] (0mm,3.0mm)--(-3.0mm,-2.2mm)--(3.0mm,-2.2mm)--cycle;
\node[font={\sffamily\fontsize{8.888888888888888}{8.888888888888888}\selectfont\bfseries},text=white,inner sep=0pt] at (0mm,-.35mm) {2};
\end{scope}
\node[g2tree dark,circle,minimum size=6.48mm,font={\sffamily\fontsize{12}{12}\selectfont\bfseries}] at (v3) {3};
\node[g2tree light,rectangle,minimum size=6.21mm,font={\sffamily\fontsize{12}{12}\selectfont\bfseries}] at (v4) {4};
\begin{scope}[shift={(v5)},scale=1.35,transform shape]
\path[fill=g2TaxonYellow,draw=white,line width=.55pt,line join=round] (0mm,3.0mm)--(-3.0mm,-2.2mm)--(3.0mm,-2.2mm)--cycle;
\node[font={\sffamily\fontsize{8.888888888888888}{8.888888888888888}\selectfont\bfseries},text=white,inner sep=0pt] at (0mm,-.35mm) {5};
\end{scope}
\node[g2tree light,circle,minimum size=6.48mm,font={\sffamily\fontsize{12}{12}\selectfont\bfseries}] at (v6) {6};
\begin{scope}[shift={(v7)},scale=1.35,transform shape]
\path[fill=g2HeartRed,draw=white,line width=.55pt]
  (0mm,-2.7mm)..controls(-.45mm,-2.1mm)and(-3.0mm,.35mm)..(-3.0mm,1.45mm)
  ..controls(-3.0mm,3.2mm)and(-.9mm,3.4mm)..(0mm,1.75mm)
  ..controls(.9mm,3.4mm)and(3.0mm,3.2mm)..(3.0mm,1.45mm)
  ..controls(3.0mm,.35mm)and(.45mm,-2.1mm)..(0mm,-2.7mm)--cycle;
\node[font={\sffamily\fontsize{8.888888888888888}{8.888888888888888}\selectfont\bfseries},text=white,inner sep=0pt] at (0mm,.35mm) {7};
\end{scope}
\end{tikzpicture}
\endgroup
}};
  \node[ft tree caption] at (2.65,16.12) {representative cell 70\\rays $7,10,33$};
  \node[font=\fontsize{14}{16}\selectfont\bfseries,text=ftQQBlue,align=center] at (7.1,19.85)
    {$ABC$ and $ABD$};
  \FTTriangle{6.10}{18.48}{ftQQBlue}{A}{B}{C}
  \FTTriangle{8.02}{18.48}{ftQQBlue}{A}{B}{D}
  \node[ft small,text width=3.8cm] at (7.1,16.90)
    {2 faces / tetrahedron\\cells 159 and 71};

  \path[rounded corners=3pt,fill=ftQQTeal!4,draw=ftQQTeal!38,line width=.8pt]
    (0.05,8.65) rectangle (9.25,14.95);
  \node[ft title] at (4.65,14.28) {$\mathbb Q$ orbit 28\qquad size 3};
  \node at (2.65,11.68) {\resizebox{4.25cm}{!}{% Generated by export_field_atlas_assets.jl from the maintained census.
\begingroup
\definecolor{g2TaxonBlue}{HTML}{1676A3}
\definecolor{g2TaxonYellow}{HTML}{E39A24}
\definecolor{g2HeartRed}{HTML}{D7443E}
\definecolor{g2EdgeBlue}{HTML}{205B7A}
\tikzset{
  g2tree pendant/.style={draw=black!43,line width=.6pt,line cap=round},
  g2tree internal/.style={draw=g2EdgeBlue,line width=2pt,line cap=round},
  g2tree vertex/.style={circle,fill=g2EdgeBlue,draw=white,line width=.4pt,inner sep=0pt,minimum size=4.8pt},
  g2tree dark/.style={fill=g2TaxonBlue,draw=white,line width=.55pt,text=white,inner sep=0pt},
  g2tree light/.style={fill=g2TaxonYellow,draw=white,line width=.55pt,text=white,inner sep=0pt},
  g2tree weight/.style={circle,fill=white,draw=none,text=black,inner sep=.15pt,minimum size=10.0pt},
  g2tree aligned pendant/.style={fill=white,text=g2EdgeBlue,inner sep=.7pt,sloped,above=1pt,font={\sffamily\fontsize{9.25}{9.75}\selectfont}},
  g2tree aligned internal/.style={fill=white,text=g2EdgeBlue,inner sep=1pt,sloped,above=1pt,font={\sffamily\fontsize{9.25}{9.75}\selectfont\bfseries}},
}
\begin{tikzpicture}[x=1.0cm,y=1.0cm,baseline=(current bounding box.center)]
\coordinate (v1) at (1.65748,1.32180);
\coordinate (v2) at (0.91983,-1.91005);
\coordinate (v3) at (-2.06685,-0.47174);
\coordinate (v4) at (-0.91983,-1.91005);
\coordinate (v5) at (2.06685,-0.47174);
\coordinate (v6) at (-1.65748,1.32180);
\coordinate (v7) at (0.00000,2.12000);
\coordinate (v8) at (0.65172,0.22413);
\coordinate (v9) at (-0.07349,-0.68525);
\coordinate (v10) at (-0.01547,0.01234);
\coordinate (v11) at (-0.56276,0.44878);
\draw[g2tree internal] (v8)--(v10);
\draw[g2tree internal] (v9)--(v10);
\draw[g2tree internal] (v10)--(v11);
\draw[g2tree pendant] (v1)--(v8);
\draw[g2tree pendant] (v2)--(v9);
\draw[g2tree pendant] (v3)--(v11);
\draw[g2tree pendant] (v4)--(v9);
\draw[g2tree pendant] (v5)--(v8);
\draw[g2tree pendant] (v6)--(v11);
\draw[g2tree pendant] (v7)--(v11);
\node[g2tree weight,font={\sffamily\fontsize{9.25}{9.75}\selectfont}] at (0.31812,0.11823) {2};
\node[g2tree weight,font={\sffamily\fontsize{9.25}{9.75}\selectfont}] at (-0.04448,-0.33646) {2};
\node[g2tree weight,font={\sffamily\fontsize{9.25}{9.75}\selectfont}] at (-0.28911,0.23056) {1};
\node[g2tree weight,font={\sffamily\fontsize{9.25}{9.75}\selectfont}] at (1.15460,0.77296) {0};
\node[g2tree weight,font={\sffamily\fontsize{9.25}{9.75}\selectfont}] at (0.42317,-1.29765) {0};
\node[g2tree weight,font={\sffamily\fontsize{9.25}{9.75}\selectfont}] at (-1.31480,-0.01148) {2};
\node[g2tree weight,font={\sffamily\fontsize{9.25}{9.75}\selectfont}] at (-0.49666,-1.29765) {0};
\node[g2tree weight,font={\sffamily\fontsize{9.25}{9.75}\selectfont}] at (1.35928,-0.12381) {0};
\node[g2tree weight,font={\sffamily\fontsize{9.25}{9.75}\selectfont}] at (-1.11012,0.88529) {2};
\node[g2tree weight,font={\sffamily\fontsize{9.25}{9.75}\selectfont}] at (-0.28138,1.28439) {2};
\node[g2tree vertex] at (v8) {};
\node[g2tree vertex] at (v9) {};
\node[g2tree vertex] at (v10) {};
\node[g2tree vertex] at (v11) {};
\node[g2tree dark,rectangle,minimum size=6.21mm,font={\sffamily\fontsize{12}{12}\selectfont\bfseries}] at (v1) {1};
\begin{scope}[shift={(v2)},scale=1.35,transform shape]
\path[fill=g2TaxonBlue,draw=white,line width=.55pt,line join=round] (0mm,3.0mm)--(-3.0mm,-2.2mm)--(3.0mm,-2.2mm)--cycle;
\node[font={\sffamily\fontsize{8.888888888888888}{8.888888888888888}\selectfont\bfseries},text=white,inner sep=0pt] at (0mm,-.35mm) {2};
\end{scope}
\node[g2tree dark,circle,minimum size=6.48mm,font={\sffamily\fontsize{12}{12}\selectfont\bfseries}] at (v3) {3};
\node[g2tree light,rectangle,minimum size=6.21mm,font={\sffamily\fontsize{12}{12}\selectfont\bfseries}] at (v4) {4};
\begin{scope}[shift={(v5)},scale=1.35,transform shape]
\path[fill=g2TaxonYellow,draw=white,line width=.55pt,line join=round] (0mm,3.0mm)--(-3.0mm,-2.2mm)--(3.0mm,-2.2mm)--cycle;
\node[font={\sffamily\fontsize{8.888888888888888}{8.888888888888888}\selectfont\bfseries},text=white,inner sep=0pt] at (0mm,-.35mm) {5};
\end{scope}
\node[g2tree light,circle,minimum size=6.48mm,font={\sffamily\fontsize{12}{12}\selectfont\bfseries}] at (v6) {6};
\begin{scope}[shift={(v7)},scale=1.35,transform shape]
\path[fill=g2HeartRed,draw=white,line width=.55pt]
  (0mm,-2.7mm)..controls(-.45mm,-2.1mm)and(-3.0mm,.35mm)..(-3.0mm,1.45mm)
  ..controls(-3.0mm,3.2mm)and(-.9mm,3.4mm)..(0mm,1.75mm)
  ..controls(.9mm,3.4mm)and(3.0mm,3.2mm)..(3.0mm,1.45mm)
  ..controls(3.0mm,.35mm)and(.45mm,-2.1mm)..(0mm,-2.7mm)--cycle;
\node[font={\sffamily\fontsize{8.888888888888888}{8.888888888888888}\selectfont\bfseries},text=white,inner sep=0pt] at (0mm,.35mm) {7};
\end{scope}
\end{tikzpicture}
\endgroup
}};
  \node[ft tree caption] at (2.65,9.22) {representative cell 120\\rays $8,11,12$};
  \node[font=\fontsize{14}{16}\selectfont\bfseries,text=ftQQTeal] at (7.1,12.95) {$ACD$};
  \FTTriangle{7.1}{11.58}{ftQQTeal}{A}{C}{D}
  \node[ft small,text width=3.8cm] at (7.1,10.0)
    {1 face per tetrahedron\\displayed cell 140};

  \path[rounded corners=3pt,fill=ftQQPurple!4,draw=ftQQPurple!38,line width=.8pt]
    (0.05,1.75) rectangle (9.25,8.05);
  \node[ft title] at (4.65,7.38) {$\mathbb Q$ orbit 30\qquad size 3};
  \node at (2.65,4.78) {\resizebox{4.25cm}{!}{% Generated by export_field_atlas_assets.jl from the maintained census.
\begingroup
\definecolor{g2TaxonBlue}{HTML}{1676A3}
\definecolor{g2TaxonYellow}{HTML}{E39A24}
\definecolor{g2HeartRed}{HTML}{D7443E}
\definecolor{g2EdgeBlue}{HTML}{205B7A}
\tikzset{
  g2tree pendant/.style={draw=black!43,line width=.6pt,line cap=round},
  g2tree internal/.style={draw=g2EdgeBlue,line width=2pt,line cap=round},
  g2tree vertex/.style={circle,fill=g2EdgeBlue,draw=white,line width=.4pt,inner sep=0pt,minimum size=4.8pt},
  g2tree dark/.style={fill=g2TaxonBlue,draw=white,line width=.55pt,text=white,inner sep=0pt},
  g2tree light/.style={fill=g2TaxonYellow,draw=white,line width=.55pt,text=white,inner sep=0pt},
  g2tree weight/.style={circle,fill=white,draw=none,text=black,inner sep=.15pt,minimum size=10.0pt},
  g2tree aligned pendant/.style={fill=white,text=g2EdgeBlue,inner sep=.7pt,sloped,above=1pt,font={\sffamily\fontsize{9.25}{9.75}\selectfont}},
  g2tree aligned internal/.style={fill=white,text=g2EdgeBlue,inner sep=1pt,sloped,above=1pt,font={\sffamily\fontsize{9.25}{9.75}\selectfont\bfseries}},
}
\begin{tikzpicture}[x=1.0cm,y=1.0cm,baseline=(current bounding box.center)]
\coordinate (v1) at (1.65748,1.32180);
\coordinate (v2) at (0.91983,-1.91005);
\coordinate (v3) at (-2.06685,-0.47174);
\coordinate (v4) at (2.06685,-0.47174);
\coordinate (v5) at (-1.65748,1.32180);
\coordinate (v6) at (-0.91983,-1.91005);
\coordinate (v7) at (0.00000,2.12000);
\coordinate (v8) at (0.00000,-0.51634);
\coordinate (v9) at (-0.69979,0.16634);
\coordinate (v10) at (0.69979,0.16634);
\coordinate (v11) at (0.00000,0.18366);
\draw[g2tree internal] (v8)--(v11);
\draw[g2tree internal] (v9)--(v11);
\draw[g2tree internal] (v10)--(v11);
\draw[g2tree pendant] (v1)--(v10);
\draw[g2tree pendant] (v2)--(v8);
\draw[g2tree pendant] (v3)--(v9);
\draw[g2tree pendant] (v4)--(v10);
\draw[g2tree pendant] (v5)--(v9);
\draw[g2tree pendant] (v6)--(v8);
\draw[g2tree pendant] (v7)--(v11);
\node[g2tree weight,font={\sffamily\fontsize{9.25}{9.75}\selectfont}] at (0.00000,-0.16634) {2};
\node[g2tree weight,font={\sffamily\fontsize{9.25}{9.75}\selectfont}] at (-0.34989,0.17500) {2};
\node[g2tree weight,font={\sffamily\fontsize{9.25}{9.75}\selectfont}] at (0.34989,0.17500) {1};
\node[g2tree weight,font={\sffamily\fontsize{9.25}{9.75}\selectfont}] at (1.17863,0.74407) {1};
\node[g2tree weight,font={\sffamily\fontsize{9.25}{9.75}\selectfont}] at (0.45992,-1.21320) {0};
\node[g2tree weight,font={\sffamily\fontsize{9.25}{9.75}\selectfont}] at (-1.38332,-0.15270) {0};
\node[g2tree weight,font={\sffamily\fontsize{9.25}{9.75}\selectfont}] at (1.38332,-0.15270) {1};
\node[g2tree weight,font={\sffamily\fontsize{9.25}{9.75}\selectfont}] at (-1.17863,0.74407) {0};
\node[g2tree weight,font={\sffamily\fontsize{9.25}{9.75}\selectfont}] at (-0.45992,-1.21320) {0};
\node[g2tree weight,font={\sffamily\fontsize{9.25}{9.75}\selectfont}] at (0.00000,1.15183) {2};
\node[g2tree vertex] at (v8) {};
\node[g2tree vertex] at (v9) {};
\node[g2tree vertex] at (v10) {};
\node[g2tree vertex] at (v11) {};
\node[g2tree dark,rectangle,minimum size=6.21mm,font={\sffamily\fontsize{12}{12}\selectfont\bfseries}] at (v1) {1};
\begin{scope}[shift={(v2)},scale=1.35,transform shape]
\path[fill=g2TaxonBlue,draw=white,line width=.55pt,line join=round] (0mm,3.0mm)--(-3.0mm,-2.2mm)--(3.0mm,-2.2mm)--cycle;
\node[font={\sffamily\fontsize{8.888888888888888}{8.888888888888888}\selectfont\bfseries},text=white,inner sep=0pt] at (0mm,-.35mm) {2};
\end{scope}
\node[g2tree dark,circle,minimum size=6.48mm,font={\sffamily\fontsize{12}{12}\selectfont\bfseries}] at (v3) {3};
\node[g2tree light,rectangle,minimum size=6.21mm,font={\sffamily\fontsize{12}{12}\selectfont\bfseries}] at (v4) {4};
\begin{scope}[shift={(v5)},scale=1.35,transform shape]
\path[fill=g2TaxonYellow,draw=white,line width=.55pt,line join=round] (0mm,3.0mm)--(-3.0mm,-2.2mm)--(3.0mm,-2.2mm)--cycle;
\node[font={\sffamily\fontsize{8.888888888888888}{8.888888888888888}\selectfont\bfseries},text=white,inner sep=0pt] at (0mm,-.35mm) {5};
\end{scope}
\node[g2tree light,circle,minimum size=6.48mm,font={\sffamily\fontsize{12}{12}\selectfont\bfseries}] at (v6) {6};
\begin{scope}[shift={(v7)},scale=1.35,transform shape]
\path[fill=g2HeartRed,draw=white,line width=.55pt]
  (0mm,-2.7mm)..controls(-.45mm,-2.1mm)and(-3.0mm,.35mm)..(-3.0mm,1.45mm)
  ..controls(-3.0mm,3.2mm)and(-.9mm,3.4mm)..(0mm,1.75mm)
  ..controls(.9mm,3.4mm)and(3.0mm,3.2mm)..(3.0mm,1.45mm)
  ..controls(3.0mm,.35mm)and(.45mm,-2.1mm)..(0mm,-2.7mm)--cycle;
\node[font={\sffamily\fontsize{8.888888888888888}{8.888888888888888}\selectfont\bfseries},text=white,inner sep=0pt] at (0mm,.35mm) {7};
\end{scope}
\end{tikzpicture}
\endgroup
}};
  \node[ft tree caption] at (2.65,2.32) {representative cell 132\\rays $6,9,40$};
  \node[font=\fontsize{14}{16}\selectfont\bfseries,text=ftQQPurple] at (7.1,6.05) {$BCD$};
  \FTTriangle{7.1}{4.68}{ftQQPurple}{B}{C}{D}
  \node[ft small,text width=3.8cm] at (7.1,3.1)
    {1 face per tetrahedron\\displayed cell 132};

  % Right: spatial order follows the tetrahedron, rather than orbit number.
  \path[rounded corners=3pt,fill=ftGFGold!4,draw=ftGFGold!38,line width=.8pt]
    (30.2,16.35) rectangle (39.4,21.25);
  \node[ft title] at (34.8,20.68) {$\operatorname{GF}(2)$ orbit 20\qquad size 3};
  \node at (32.5,18.66) {\resizebox{3.65cm}{!}{% Generated by export_field_atlas_assets.jl from the maintained census.
\begingroup
\definecolor{g2TaxonBlue}{HTML}{1676A3}
\definecolor{g2TaxonYellow}{HTML}{E39A24}
\definecolor{g2HeartRed}{HTML}{D7443E}
\definecolor{g2EdgeBlue}{HTML}{205B7A}
\tikzset{
  g2tree pendant/.style={draw=black!43,line width=.6pt,line cap=round},
  g2tree internal/.style={draw=g2EdgeBlue,line width=2pt,line cap=round},
  g2tree vertex/.style={circle,fill=g2EdgeBlue,draw=white,line width=.4pt,inner sep=0pt,minimum size=4.8pt},
  g2tree dark/.style={fill=g2TaxonBlue,draw=white,line width=.55pt,text=white,inner sep=0pt},
  g2tree light/.style={fill=g2TaxonYellow,draw=white,line width=.55pt,text=white,inner sep=0pt},
  g2tree weight/.style={circle,fill=white,draw=none,text=black,inner sep=.15pt,minimum size=10.0pt},
  g2tree aligned pendant/.style={fill=white,text=g2EdgeBlue,inner sep=.7pt,sloped,above=1pt,font={\sffamily\fontsize{9.25}{9.75}\selectfont}},
  g2tree aligned internal/.style={fill=white,text=g2EdgeBlue,inner sep=1pt,sloped,above=1pt,font={\sffamily\fontsize{9.25}{9.75}\selectfont\bfseries}},
}
\begin{tikzpicture}[x=1.0cm,y=1.0cm,baseline=(current bounding box.center)]
\coordinate (v1) at (1.65748,1.32180);
\coordinate (v2) at (0.91983,-1.91005);
\coordinate (v3) at (-2.06685,-0.47174);
\coordinate (v4) at (2.06685,-0.47174);
\coordinate (v5) at (-1.65748,1.32180);
\coordinate (v6) at (-0.91983,-1.91005);
\coordinate (v7) at (0.00000,2.12000);
\coordinate (v8) at (-0.05494,-0.77281);
\coordinate (v9) at (-0.86696,0.12576);
\coordinate (v10) at (0.89666,0.26787);
\coordinate (v11) at (0.22566,0.46729);
\coordinate (v12) at (-0.20042,-0.08810);
\draw[g2tree internal] (v8)--(v12);
\draw[g2tree internal] (v9)--(v12);
\draw[g2tree internal] (v10)--(v11);
\draw[g2tree internal] (v11)--(v12);
\draw[g2tree pendant] (v1)--(v10);
\draw[g2tree pendant] (v2)--(v8);
\draw[g2tree pendant] (v3)--(v9);
\draw[g2tree pendant] (v4)--(v10);
\draw[g2tree pendant] (v5)--(v9);
\draw[g2tree pendant] (v6)--(v8);
\draw[g2tree pendant] (v7)--(v11);
\node[g2tree weight,font={\sffamily\fontsize{9.25}{9.75}\selectfont}] at (-0.12768,-0.43046) {3};
\node[g2tree weight,font={\sffamily\fontsize{9.25}{9.75}\selectfont}] at (-0.53369,0.01883) {3};
\node[g2tree weight,font={\sffamily\fontsize{9.25}{9.75}\selectfont}] at (0.56116,0.36758) {1};
\node[g2tree weight,font={\sffamily\fontsize{9.25}{9.75}\selectfont}] at (0.01262,0.18959) {1};
\node[g2tree weight,font={\sffamily\fontsize{9.25}{9.75}\selectfont}] at (1.27707,0.79483) {2};
\node[g2tree weight,font={\sffamily\fontsize{9.25}{9.75}\selectfont}] at (0.43245,-1.34143) {0};
\node[g2tree weight,font={\sffamily\fontsize{9.25}{9.75}\selectfont}] at (-1.46690,-0.17299) {0};
\node[g2tree weight,font={\sffamily\fontsize{9.25}{9.75}\selectfont}] at (1.48175,-0.10194) {2};
\node[g2tree weight,font={\sffamily\fontsize{9.25}{9.75}\selectfont}] at (-1.26222,0.72378) {0};
\node[g2tree weight,font={\sffamily\fontsize{9.25}{9.75}\selectfont}] at (-0.48739,-1.34143) {0};
\node[g2tree weight,font={\sffamily\fontsize{9.25}{9.75}\selectfont}] at (0.11283,1.29364) {3};
\node[g2tree vertex] at (v8) {};
\node[g2tree vertex] at (v9) {};
\node[g2tree vertex] at (v10) {};
\node[g2tree vertex] at (v11) {};
\node[g2tree vertex] at (v12) {};
\node[g2tree dark,rectangle,minimum size=6.21mm,font={\sffamily\fontsize{12}{12}\selectfont\bfseries}] at (v1) {1};
\begin{scope}[shift={(v2)},scale=1.35,transform shape]
\path[fill=g2TaxonBlue,draw=white,line width=.55pt,line join=round] (0mm,3.0mm)--(-3.0mm,-2.2mm)--(3.0mm,-2.2mm)--cycle;
\node[font={\sffamily\fontsize{8.888888888888888}{8.888888888888888}\selectfont\bfseries},text=white,inner sep=0pt] at (0mm,-.35mm) {2};
\end{scope}
\node[g2tree dark,circle,minimum size=6.48mm,font={\sffamily\fontsize{12}{12}\selectfont\bfseries}] at (v3) {3};
\node[g2tree light,rectangle,minimum size=6.21mm,font={\sffamily\fontsize{12}{12}\selectfont\bfseries}] at (v4) {4};
\begin{scope}[shift={(v5)},scale=1.35,transform shape]
\path[fill=g2TaxonYellow,draw=white,line width=.55pt,line join=round] (0mm,3.0mm)--(-3.0mm,-2.2mm)--(3.0mm,-2.2mm)--cycle;
\node[font={\sffamily\fontsize{8.888888888888888}{8.888888888888888}\selectfont\bfseries},text=white,inner sep=0pt] at (0mm,-.35mm) {5};
\end{scope}
\node[g2tree light,circle,minimum size=6.48mm,font={\sffamily\fontsize{12}{12}\selectfont\bfseries}] at (v6) {6};
\begin{scope}[shift={(v7)},scale=1.35,transform shape]
\path[fill=g2HeartRed,draw=white,line width=.55pt]
  (0mm,-2.7mm)..controls(-.45mm,-2.1mm)and(-3.0mm,.35mm)..(-3.0mm,1.45mm)
  ..controls(-3.0mm,3.2mm)and(-.9mm,3.4mm)..(0mm,1.75mm)
  ..controls(.9mm,3.4mm)and(3.0mm,3.2mm)..(3.0mm,1.45mm)
  ..controls(3.0mm,.35mm)and(.45mm,-2.1mm)..(0mm,-2.7mm)--cycle;
\node[font={\sffamily\fontsize{8.888888888888888}{8.888888888888888}\selectfont\bfseries},text=white,inner sep=0pt] at (0mm,.35mm) {7};
\end{scope}
\end{tikzpicture}
\endgroup
}};
  \node[ft tree caption] at (32.5,16.65) {rep. cell 37\quad rays $3,13,14$};
  \node[font=\fontsize{14}{16}\selectfont\bfseries,text=ftGFGold] at (36.95,19.64) {$OCD$};
  \FTTriangle{36.95}{18.36}{ftGFGold}{C}{D}{O}
  \node[ft small] at (36.95,17.07) {1 triangle per tetrahedron};

  \path[rounded corners=3pt,fill=ftGFTeal!4,draw=ftGFTeal!38,line width=.8pt]
    (30.2,11.05) rectangle (39.4,15.95);
  \node[ft title] at (34.8,15.38) {$\operatorname{GF}(2)$ orbit 9\qquad size 6};
  \node at (32.5,13.36) {\resizebox{3.65cm}{!}{% Generated by export_field_atlas_assets.jl from the maintained census.
\begingroup
\definecolor{g2TaxonBlue}{HTML}{1676A3}
\definecolor{g2TaxonYellow}{HTML}{E39A24}
\definecolor{g2HeartRed}{HTML}{D7443E}
\definecolor{g2EdgeBlue}{HTML}{205B7A}
\tikzset{
  g2tree pendant/.style={draw=black!43,line width=.6pt,line cap=round},
  g2tree internal/.style={draw=g2EdgeBlue,line width=2pt,line cap=round},
  g2tree vertex/.style={circle,fill=g2EdgeBlue,draw=white,line width=.4pt,inner sep=0pt,minimum size=4.8pt},
  g2tree dark/.style={fill=g2TaxonBlue,draw=white,line width=.55pt,text=white,inner sep=0pt},
  g2tree light/.style={fill=g2TaxonYellow,draw=white,line width=.55pt,text=white,inner sep=0pt},
  g2tree weight/.style={circle,fill=white,draw=none,text=black,inner sep=.15pt,minimum size=10.0pt},
  g2tree aligned pendant/.style={fill=white,text=g2EdgeBlue,inner sep=.7pt,sloped,above=1pt,font={\sffamily\fontsize{9.25}{9.75}\selectfont}},
  g2tree aligned internal/.style={fill=white,text=g2EdgeBlue,inner sep=1pt,sloped,above=1pt,font={\sffamily\fontsize{9.25}{9.75}\selectfont\bfseries}},
}
\begin{tikzpicture}[x=1.0cm,y=1.0cm,baseline=(current bounding box.center)]
\coordinate (v1) at (1.65748,1.32180);
\coordinate (v2) at (0.91983,-1.91005);
\coordinate (v3) at (-2.06685,-0.47174);
\coordinate (v4) at (2.06685,-0.47174);
\coordinate (v5) at (-1.65748,1.32180);
\coordinate (v6) at (-0.91983,-1.91005);
\coordinate (v7) at (0.00000,2.12000);
\coordinate (v8) at (-0.05494,-0.77281);
\coordinate (v9) at (0.89666,0.26787);
\coordinate (v10) at (-0.86696,0.12576);
\coordinate (v11) at (0.22566,0.46729);
\coordinate (v12) at (-0.20042,-0.08810);
\draw[g2tree internal] (v8)--(v12);
\draw[g2tree internal] (v9)--(v11);
\draw[g2tree internal] (v10)--(v12);
\draw[g2tree internal] (v11)--(v12);
\draw[g2tree pendant] (v1)--(v9);
\draw[g2tree pendant] (v2)--(v8);
\draw[g2tree pendant] (v3)--(v10);
\draw[g2tree pendant] (v4)--(v9);
\draw[g2tree pendant] (v5)--(v10);
\draw[g2tree pendant] (v6)--(v8);
\draw[g2tree pendant] (v7)--(v11);
\node[g2tree weight,font={\sffamily\fontsize{9.25}{9.75}\selectfont}] at (-0.12768,-0.43046) {3};
\node[g2tree weight,font={\sffamily\fontsize{9.25}{9.75}\selectfont}] at (0.56116,0.36758) {2};
\node[g2tree weight,font={\sffamily\fontsize{9.25}{9.75}\selectfont}] at (-0.53369,0.01883) {1};
\node[g2tree weight,font={\sffamily\fontsize{9.25}{9.75}\selectfont}] at (0.01262,0.18959) {1};
\node[g2tree weight,font={\sffamily\fontsize{9.25}{9.75}\selectfont}] at (1.27707,0.79483) {0};
\node[g2tree weight,font={\sffamily\fontsize{9.25}{9.75}\selectfont}] at (0.43245,-1.34143) {0};
\node[g2tree weight,font={\sffamily\fontsize{9.25}{9.75}\selectfont}] at (-1.46690,-0.17299) {0};
\node[g2tree weight,font={\sffamily\fontsize{9.25}{9.75}\selectfont}] at (1.48175,-0.10194) {0};
\node[g2tree weight,font={\sffamily\fontsize{9.25}{9.75}\selectfont}] at (-1.26222,0.72378) {0};
\node[g2tree weight,font={\sffamily\fontsize{9.25}{9.75}\selectfont}] at (-0.48739,-1.34143) {0};
\node[g2tree weight,font={\sffamily\fontsize{9.25}{9.75}\selectfont}] at (0.11283,1.29364) {2};
\node[g2tree vertex] at (v8) {};
\node[g2tree vertex] at (v9) {};
\node[g2tree vertex] at (v10) {};
\node[g2tree vertex] at (v11) {};
\node[g2tree vertex] at (v12) {};
\node[g2tree dark,rectangle,minimum size=6.21mm,font={\sffamily\fontsize{12}{12}\selectfont\bfseries}] at (v1) {1};
\begin{scope}[shift={(v2)},scale=1.35,transform shape]
\path[fill=g2TaxonBlue,draw=white,line width=.55pt,line join=round] (0mm,3.0mm)--(-3.0mm,-2.2mm)--(3.0mm,-2.2mm)--cycle;
\node[font={\sffamily\fontsize{8.888888888888888}{8.888888888888888}\selectfont\bfseries},text=white,inner sep=0pt] at (0mm,-.35mm) {2};
\end{scope}
\node[g2tree dark,circle,minimum size=6.48mm,font={\sffamily\fontsize{12}{12}\selectfont\bfseries}] at (v3) {3};
\node[g2tree light,rectangle,minimum size=6.21mm,font={\sffamily\fontsize{12}{12}\selectfont\bfseries}] at (v4) {4};
\begin{scope}[shift={(v5)},scale=1.35,transform shape]
\path[fill=g2TaxonYellow,draw=white,line width=.55pt,line join=round] (0mm,3.0mm)--(-3.0mm,-2.2mm)--(3.0mm,-2.2mm)--cycle;
\node[font={\sffamily\fontsize{8.888888888888888}{8.888888888888888}\selectfont\bfseries},text=white,inner sep=0pt] at (0mm,-.35mm) {5};
\end{scope}
\node[g2tree light,circle,minimum size=6.48mm,font={\sffamily\fontsize{12}{12}\selectfont\bfseries}] at (v6) {6};
\begin{scope}[shift={(v7)},scale=1.35,transform shape]
\path[fill=g2HeartRed,draw=white,line width=.55pt]
  (0mm,-2.7mm)..controls(-.45mm,-2.1mm)and(-3.0mm,.35mm)..(-3.0mm,1.45mm)
  ..controls(-3.0mm,3.2mm)and(-.9mm,3.4mm)..(0mm,1.75mm)
  ..controls(.9mm,3.4mm)and(3.0mm,3.2mm)..(3.0mm,1.45mm)
  ..controls(3.0mm,.35mm)and(.45mm,-2.1mm)..(0mm,-2.7mm)--cycle;
\node[font={\sffamily\fontsize{8.888888888888888}{8.888888888888888}\selectfont\bfseries},text=white,inner sep=0pt] at (0mm,.35mm) {7};
\end{scope}
\end{tikzpicture}
\endgroup
}};
  \node[ft tree caption] at (32.5,11.35) {rep. cell 10\quad rays $2,3,13$};
  \node[font=\fontsize{14}{16}\selectfont\bfseries,text=ftGFTeal] at (36.95,14.34) {$OBC$ and $OBD$};
  \FTTriangle{35.98}{13.06}{ftGFTeal}{B}{C}{O}
  \FTTriangle{37.92}{13.06}{ftGFTeal}{B}{D}{O}
  \node[ft small] at (36.95,11.77) {2 triangles per tetrahedron};

  \path[rounded corners=3pt,fill=ftGFPurple!4,draw=ftGFPurple!38,line width=.8pt]
    (30.2,5.75) rectangle (39.4,10.65);
  \node[ft title] at (34.8,10.08) {$\operatorname{GF}(2)$ orbit 10\qquad size 6};
  \node at (32.5,8.06) {\resizebox{3.65cm}{!}{% Generated by export_field_atlas_assets.jl from the maintained census.
\begingroup
\definecolor{g2TaxonBlue}{HTML}{1676A3}
\definecolor{g2TaxonYellow}{HTML}{E39A24}
\definecolor{g2HeartRed}{HTML}{D7443E}
\definecolor{g2EdgeBlue}{HTML}{205B7A}
\tikzset{
  g2tree pendant/.style={draw=black!43,line width=.6pt,line cap=round},
  g2tree internal/.style={draw=g2EdgeBlue,line width=2pt,line cap=round},
  g2tree vertex/.style={circle,fill=g2EdgeBlue,draw=white,line width=.4pt,inner sep=0pt,minimum size=4.8pt},
  g2tree dark/.style={fill=g2TaxonBlue,draw=white,line width=.55pt,text=white,inner sep=0pt},
  g2tree light/.style={fill=g2TaxonYellow,draw=white,line width=.55pt,text=white,inner sep=0pt},
  g2tree weight/.style={circle,fill=white,draw=none,text=black,inner sep=.15pt,minimum size=10.0pt},
  g2tree aligned pendant/.style={fill=white,text=g2EdgeBlue,inner sep=.7pt,sloped,above=1pt,font={\sffamily\fontsize{9.25}{9.75}\selectfont}},
  g2tree aligned internal/.style={fill=white,text=g2EdgeBlue,inner sep=1pt,sloped,above=1pt,font={\sffamily\fontsize{9.25}{9.75}\selectfont\bfseries}},
}
\begin{tikzpicture}[x=1.0cm,y=1.0cm,baseline=(current bounding box.center)]
\coordinate (v1) at (1.65748,1.32180);
\coordinate (v2) at (0.91983,-1.91005);
\coordinate (v3) at (-2.06685,-0.47174);
\coordinate (v4) at (2.06685,-0.47174);
\coordinate (v5) at (-1.65748,1.32180);
\coordinate (v6) at (-0.91983,-1.91005);
\coordinate (v7) at (0.00000,2.12000);
\coordinate (v8) at (-0.05494,-0.77281);
\coordinate (v9) at (0.89666,0.26787);
\coordinate (v10) at (0.22566,0.46729);
\coordinate (v11) at (-0.86696,0.12576);
\coordinate (v12) at (-0.20042,-0.08810);
\draw[g2tree internal] (v8)--(v12);
\draw[g2tree internal] (v9)--(v10);
\draw[g2tree internal] (v10)--(v12);
\draw[g2tree internal] (v11)--(v12);
\draw[g2tree pendant] (v1)--(v9);
\draw[g2tree pendant] (v2)--(v8);
\draw[g2tree pendant] (v3)--(v11);
\draw[g2tree pendant] (v4)--(v9);
\draw[g2tree pendant] (v5)--(v11);
\draw[g2tree pendant] (v6)--(v8);
\draw[g2tree pendant] (v7)--(v10);
\node[g2tree weight,font={\sffamily\fontsize{9.25}{9.75}\selectfont}] at (-0.12768,-0.43046) {3};
\node[g2tree weight,font={\sffamily\fontsize{9.25}{9.75}\selectfont}] at (0.56116,0.36758) {1};
\node[g2tree weight,font={\sffamily\fontsize{9.25}{9.75}\selectfont}] at (0.01262,0.18959) {2};
\node[g2tree weight,font={\sffamily\fontsize{9.25}{9.75}\selectfont}] at (-0.53369,0.01883) {1};
\node[g2tree weight,font={\sffamily\fontsize{9.25}{9.75}\selectfont}] at (1.27707,0.79483) {1};
\node[g2tree weight,font={\sffamily\fontsize{9.25}{9.75}\selectfont}] at (0.43245,-1.34143) {0};
\node[g2tree weight,font={\sffamily\fontsize{9.25}{9.75}\selectfont}] at (-1.46690,-0.17299) {0};
\node[g2tree weight,font={\sffamily\fontsize{9.25}{9.75}\selectfont}] at (1.48175,-0.10194) {1};
\node[g2tree weight,font={\sffamily\fontsize{9.25}{9.75}\selectfont}] at (-1.26222,0.72378) {0};
\node[g2tree weight,font={\sffamily\fontsize{9.25}{9.75}\selectfont}] at (-0.48739,-1.34143) {0};
\node[g2tree weight,font={\sffamily\fontsize{9.25}{9.75}\selectfont}] at (0.11283,1.29364) {2};
\node[g2tree vertex] at (v8) {};
\node[g2tree vertex] at (v9) {};
\node[g2tree vertex] at (v10) {};
\node[g2tree vertex] at (v11) {};
\node[g2tree vertex] at (v12) {};
\node[g2tree dark,rectangle,minimum size=6.21mm,font={\sffamily\fontsize{12}{12}\selectfont\bfseries}] at (v1) {1};
\begin{scope}[shift={(v2)},scale=1.35,transform shape]
\path[fill=g2TaxonBlue,draw=white,line width=.55pt,line join=round] (0mm,3.0mm)--(-3.0mm,-2.2mm)--(3.0mm,-2.2mm)--cycle;
\node[font={\sffamily\fontsize{8.888888888888888}{8.888888888888888}\selectfont\bfseries},text=white,inner sep=0pt] at (0mm,-.35mm) {2};
\end{scope}
\node[g2tree dark,circle,minimum size=6.48mm,font={\sffamily\fontsize{12}{12}\selectfont\bfseries}] at (v3) {3};
\node[g2tree light,rectangle,minimum size=6.21mm,font={\sffamily\fontsize{12}{12}\selectfont\bfseries}] at (v4) {4};
\begin{scope}[shift={(v5)},scale=1.35,transform shape]
\path[fill=g2TaxonYellow,draw=white,line width=.55pt,line join=round] (0mm,3.0mm)--(-3.0mm,-2.2mm)--(3.0mm,-2.2mm)--cycle;
\node[font={\sffamily\fontsize{8.888888888888888}{8.888888888888888}\selectfont\bfseries},text=white,inner sep=0pt] at (0mm,-.35mm) {5};
\end{scope}
\node[g2tree light,circle,minimum size=6.48mm,font={\sffamily\fontsize{12}{12}\selectfont\bfseries}] at (v6) {6};
\begin{scope}[shift={(v7)},scale=1.35,transform shape]
\path[fill=g2HeartRed,draw=white,line width=.55pt]
  (0mm,-2.7mm)..controls(-.45mm,-2.1mm)and(-3.0mm,.35mm)..(-3.0mm,1.45mm)
  ..controls(-3.0mm,3.2mm)and(-.9mm,3.4mm)..(0mm,1.75mm)
  ..controls(.9mm,3.4mm)and(3.0mm,3.2mm)..(3.0mm,1.45mm)
  ..controls(3.0mm,.35mm)and(.45mm,-2.1mm)..(0mm,-2.7mm)--cycle;
\node[font={\sffamily\fontsize{8.888888888888888}{8.888888888888888}\selectfont\bfseries},text=white,inner sep=0pt] at (0mm,.35mm) {7};
\end{scope}
\end{tikzpicture}
\endgroup
}};
  \node[ft tree caption] at (32.5,6.05) {rep. cell 12\quad rays $1,3,13$};
  \node[font=\fontsize{14}{16}\selectfont\bfseries,text=ftGFPurple] at (36.95,9.04) {$OAC$ and $OAD$};
  \FTTriangle{35.98}{7.76}{ftGFPurple}{A}{C}{O}
  \FTTriangle{37.92}{7.76}{ftGFPurple}{A}{D}{O}
  \node[ft small] at (36.95,6.47) {2 triangles per tetrahedron};

  \path[rounded corners=3pt,fill=ftGFRed!4,draw=ftGFRed!38,line width=.8pt]
    (30.2,.45) rectangle (39.4,5.35);
  \node[ft title] at (34.8,4.78) {$\operatorname{GF}(2)$ orbit 1\qquad size 3};
  \node at (32.5,2.76) {\resizebox{3.65cm}{!}{% Generated by export_field_atlas_assets.jl from the maintained census.
\begingroup
\definecolor{g2TaxonBlue}{HTML}{1676A3}
\definecolor{g2TaxonYellow}{HTML}{E39A24}
\definecolor{g2HeartRed}{HTML}{D7443E}
\definecolor{g2EdgeBlue}{HTML}{205B7A}
\tikzset{
  g2tree pendant/.style={draw=black!43,line width=.6pt,line cap=round},
  g2tree internal/.style={draw=g2EdgeBlue,line width=2pt,line cap=round},
  g2tree vertex/.style={circle,fill=g2EdgeBlue,draw=white,line width=.4pt,inner sep=0pt,minimum size=4.8pt},
  g2tree dark/.style={fill=g2TaxonBlue,draw=white,line width=.55pt,text=white,inner sep=0pt},
  g2tree light/.style={fill=g2TaxonYellow,draw=white,line width=.55pt,text=white,inner sep=0pt},
  g2tree weight/.style={circle,fill=white,draw=none,text=black,inner sep=.15pt,minimum size=10.0pt},
  g2tree aligned pendant/.style={fill=white,text=g2EdgeBlue,inner sep=.7pt,sloped,above=1pt,font={\sffamily\fontsize{9.25}{9.75}\selectfont}},
  g2tree aligned internal/.style={fill=white,text=g2EdgeBlue,inner sep=1pt,sloped,above=1pt,font={\sffamily\fontsize{9.25}{9.75}\selectfont\bfseries}},
}
\begin{tikzpicture}[x=1.0cm,y=1.0cm,baseline=(current bounding box.center)]
\coordinate (v1) at (1.65748,1.32180);
\coordinate (v2) at (0.91983,-1.91005);
\coordinate (v3) at (-2.06685,-0.47174);
\coordinate (v4) at (2.06685,-0.47174);
\coordinate (v5) at (-1.65748,1.32180);
\coordinate (v6) at (-0.91983,-1.91005);
\coordinate (v7) at (0.00000,2.12000);
\coordinate (v8) at (0.89666,0.26787);
\coordinate (v9) at (0.22566,0.46729);
\coordinate (v10) at (-0.05494,-0.77281);
\coordinate (v11) at (-0.86696,0.12576);
\coordinate (v12) at (-0.20042,-0.08810);
\draw[g2tree internal] (v8)--(v9);
\draw[g2tree internal] (v9)--(v12);
\draw[g2tree internal] (v10)--(v12);
\draw[g2tree internal] (v11)--(v12);
\draw[g2tree pendant] (v1)--(v8);
\draw[g2tree pendant] (v2)--(v10);
\draw[g2tree pendant] (v3)--(v11);
\draw[g2tree pendant] (v4)--(v8);
\draw[g2tree pendant] (v5)--(v11);
\draw[g2tree pendant] (v6)--(v10);
\draw[g2tree pendant] (v7)--(v9);
\node[g2tree weight,font={\sffamily\fontsize{9.25}{9.75}\selectfont}] at (0.56116,0.36758) {2};
\node[g2tree weight,font={\sffamily\fontsize{9.25}{9.75}\selectfont}] at (0.01262,0.18959) {2};
\node[g2tree weight,font={\sffamily\fontsize{9.25}{9.75}\selectfont}] at (-0.12768,-0.43046) {1};
\node[g2tree weight,font={\sffamily\fontsize{9.25}{9.75}\selectfont}] at (-0.53369,0.01883) {1};
\node[g2tree weight,font={\sffamily\fontsize{9.25}{9.75}\selectfont}] at (1.27707,0.79483) {0};
\node[g2tree weight,font={\sffamily\fontsize{9.25}{9.75}\selectfont}] at (0.43245,-1.34143) {1};
\node[g2tree weight,font={\sffamily\fontsize{9.25}{9.75}\selectfont}] at (-1.46690,-0.17299) {1};
\node[g2tree weight,font={\sffamily\fontsize{9.25}{9.75}\selectfont}] at (1.48175,-0.10194) {0};
\node[g2tree weight,font={\sffamily\fontsize{9.25}{9.75}\selectfont}] at (-1.26222,0.72378) {1};
\node[g2tree weight,font={\sffamily\fontsize{9.25}{9.75}\selectfont}] at (-0.48739,-1.34143) {1};
\node[g2tree weight,font={\sffamily\fontsize{9.25}{9.75}\selectfont}] at (0.11283,1.29364) {2};
\node[g2tree vertex] at (v8) {};
\node[g2tree vertex] at (v9) {};
\node[g2tree vertex] at (v10) {};
\node[g2tree vertex] at (v11) {};
\node[g2tree vertex] at (v12) {};
\node[g2tree dark,rectangle,minimum size=6.21mm,font={\sffamily\fontsize{12}{12}\selectfont\bfseries}] at (v1) {1};
\begin{scope}[shift={(v2)},scale=1.35,transform shape]
\path[fill=g2TaxonBlue,draw=white,line width=.55pt,line join=round] (0mm,3.0mm)--(-3.0mm,-2.2mm)--(3.0mm,-2.2mm)--cycle;
\node[font={\sffamily\fontsize{8.888888888888888}{8.888888888888888}\selectfont\bfseries},text=white,inner sep=0pt] at (0mm,-.35mm) {2};
\end{scope}
\node[g2tree dark,circle,minimum size=6.48mm,font={\sffamily\fontsize{12}{12}\selectfont\bfseries}] at (v3) {3};
\node[g2tree light,rectangle,minimum size=6.21mm,font={\sffamily\fontsize{12}{12}\selectfont\bfseries}] at (v4) {4};
\begin{scope}[shift={(v5)},scale=1.35,transform shape]
\path[fill=g2TaxonYellow,draw=white,line width=.55pt,line join=round] (0mm,3.0mm)--(-3.0mm,-2.2mm)--(3.0mm,-2.2mm)--cycle;
\node[font={\sffamily\fontsize{8.888888888888888}{8.888888888888888}\selectfont\bfseries},text=white,inner sep=0pt] at (0mm,-.35mm) {5};
\end{scope}
\node[g2tree light,circle,minimum size=6.48mm,font={\sffamily\fontsize{12}{12}\selectfont\bfseries}] at (v6) {6};
\begin{scope}[shift={(v7)},scale=1.35,transform shape]
\path[fill=g2HeartRed,draw=white,line width=.55pt]
  (0mm,-2.7mm)..controls(-.45mm,-2.1mm)and(-3.0mm,.35mm)..(-3.0mm,1.45mm)
  ..controls(-3.0mm,3.2mm)and(-.9mm,3.4mm)..(0mm,1.75mm)
  ..controls(.9mm,3.4mm)and(3.0mm,3.2mm)..(3.0mm,1.45mm)
  ..controls(3.0mm,.35mm)and(.45mm,-2.1mm)..(0mm,-2.7mm)--cycle;
\node[font={\sffamily\fontsize{8.888888888888888}{8.888888888888888}\selectfont\bfseries},text=white,inner sep=0pt] at (0mm,.35mm) {7};
\end{scope}
\end{tikzpicture}
\endgroup
}};
  \node[ft tree caption] at (32.5,.75) {rep. cell 1\quad rays $1,2,3$};
  \node[font=\fontsize{14}{16}\selectfont\bfseries,text=ftGFRed] at (36.95,3.74) {$OAB$};
  \FTTriangle{36.95}{2.46}{ftGFRed}{A}{B}{O}
  \node[ft small] at (36.95,1.17) {1 triangle per tetrahedron};

  % The two projections use exactly the same vertex positions.
  \node[font=\fontsize{16}{18}\selectfont\bfseries,text=ink] at (15.0,20.10)
    {Boundary over $\mathbb Q$};
  \node[font=\fontsize{16}{18}\selectfont\bfseries,text=ink] at (24.0,20.10)
    {Central triangles over $\operatorname{GF}(2)$};
  \begin{scope}[shift={(15.0,14.4)}]
    \coordinate (qA) at (-3.2,-2.4);
    \coordinate (qB) at (3.1,-2.7);
    \coordinate (qC) at (-1.6,4.0);
    \coordinate (qD) at (1.2,.9);
    \path[fill=ftQQBlue,fill opacity=.08] (qA)--(qB)--(qC)--cycle;
    \path[fill=ftQQBlue,fill opacity=.30] (qA)--(qB)--(qD)--cycle;
    \path[fill=ftQQTeal,fill opacity=.26] (qA)--(qC)--(qD)--cycle;
    \path[fill=ftQQPurple,fill opacity=.26] (qB)--(qC)--(qD)--cycle;
    \draw[ft edge] (qA)--(qB)--(qC)--cycle;
    \draw[ft hidden] (qA)--(qD)--(qB);
    \draw[ft edge] (qC)--(qD);
    \node[ft face label,text=ftQQBlue] (fABD) at (.15,-1.65) {$ABD$};
    \node[ft face label,text=ftQQTeal] (fACD) at (-1.72,.67) {$ACD$};
    \node[ft face label,text=ftQQPurple] (fBCD) at (.78,1.6) {$BCD$};
    \node[ft face label,text=ftQQBlue] (fABC) at (-3.38,2.8) {$ABC$};
    \draw[draw=ftQQBlue,line width=.8pt,densely dashed] (fABC.east)--(-1.40,-.25);
    \node[ft vertex] at (qA) {$A$};
    \node[ft vertex] at (qB) {$B$};
    \node[ft vertex] at (qC) {$C$};
    \node[ft vertex] at (qD) {$D$};
  \end{scope}
  \node[ft small,text width=7.6cm] at (15.0,19.07)
    {$ABC$ is the rear face.\\All four boundary faces are present.};

  \begin{scope}[shift={(24.0,14.4)}]
    \coordinate (gA) at (-3.2,-2.4);
    \coordinate (gB) at (3.1,-2.7);
    \coordinate (gC) at (-1.6,4.0);
    \coordinate (gD) at (1.2,.9);
    \coordinate (gO) at (-.125,-.05);
    \path[fill=ftGFPurple,fill opacity=.25] (gO)--(gA)--(gC)--cycle;
    \path[fill=ftGFTeal,fill opacity=.25] (gO)--(gB)--(gC)--cycle;
    \path[fill=ftGFGold,fill opacity=.34] (gO)--(gC)--(gD)--cycle;
    \path[fill=ftGFPurple,fill opacity=.21] (gO)--(gA)--(gD)--cycle;
    \path[fill=ftGFTeal,fill opacity=.21] (gO)--(gB)--(gD)--cycle;
    \path[fill=ftGFRed,fill opacity=.29] (gO)--(gA)--(gB)--cycle;
    \draw[ft edge] (gA)--(gB)--(gC)--cycle;
    \draw[ft hidden] (gA)--(gD)--(gB);
    \draw[ft edge] (gC)--(gD);
    \draw[ft centre edge] (gO)--(gA) (gO)--(gB) (gO)--(gC);
    \draw[ft centre edge,densely dashed] (gO)--(gD);
    \node[ft face label,text=ftGFGold] (fOCD) at (-.25,2.0) {$OCD$};
    \node[ft face label,text=ftGFTeal] (fOBC) at (.65,1.1) {$OBC$};
    \node[ft face label,text=ftGFTeal] (fOBD) at (1.6,-.6) {$OBD$};
    \node[ft face label,text=ftGFPurple] (fOAC) at (-1.7,.8) {$OAC$};
    \node[ft face label,text=ftGFPurple] (fOAD) at (-.85,-.62) {$OAD$};
    \node[ft face label,text=ftGFRed] (fOAB) at (-.08,-1.63) {$OAB$};
    \node[ft vertex] at (gA) {$A$};
    \node[ft vertex] at (gB) {$B$};
    \node[ft vertex] at (gC) {$C$};
    \node[ft vertex] at (gD) {$D$};
    \node[ft central] at (gO) {$O$};
  \end{scope}
  \node[ft small,text width=7.6cm] at (24.0,19.07)
    {The six edges give six triangles\\through the central ray $O$.};

  % Arrows run outward from the labelled, coloured cells to their orbit
  % trees. White underlays preserve readability when projections overlap.
  \draw[ft underlay] (fABC.west)..controls(10.7,18.5)and(10.5,18.7)..(9.40,18.7);
  \draw[ft arrow,draw=ftQQBlue] (fABC.west)..controls(10.7,18.5)and(10.5,18.7)..(9.40,18.7);
  \draw[ft underlay] (fABD.west)..controls(11.15,13.4)and(10.55,17.6)..(10.0,18.7);
  \draw[ft leader,draw=ftQQBlue] (fABD.west)..controls(11.15,13.4)and(10.55,17.6)..(10.0,18.7);
  \draw[ft underlay] (fACD.west)..controls(11.3,15.0)and(11.0,12.8)..(9.40,11.8);
  \draw[ft arrow,draw=ftQQTeal] (fACD.west)..controls(11.3,15.0)and(11.0,12.8)..(9.40,11.8);
  \draw[ft underlay] (fBCD.east)..controls(18.6,15.7)and(12.5,10.8)..(11.1,9.7)
    ..controls(10.25,8.1)and(10.8,5.8)..(9.40,4.9);
  \draw[ft arrow,draw=ftQQPurple] (fBCD.east)..controls(18.6,15.7)and(12.5,10.8)..(11.1,9.7)
    ..controls(10.25,8.1)and(10.8,5.8)..(9.40,4.9);

  \draw[ft underlay] (fOCD.east)..controls(27.1,17.3)and(28.0,18.8)..(30.05,18.8);
  \draw[ft arrow,draw=ftGFGold] (fOCD.east)..controls(27.1,17.3)and(28.0,18.8)..(30.05,18.8);
  \draw[ft underlay] (fOBC.east)..controls(27.3,16.0)and(28.0,13.5)..(30.05,13.5);
  \draw[ft arrow,draw=ftGFTeal] (fOBC.east)..controls(27.3,16.0)and(28.0,13.5)..(30.05,13.5);
  \draw[ft underlay] (fOBD.east)..controls(27.2,13.5)and(28.0,13.5)..(29.4,13.5);
  \draw[ft leader,draw=ftGFTeal] (fOBD.east)..controls(27.2,13.5)and(28.0,13.5)..(29.4,13.5);
  \draw[ft underlay] (fOAC.west)..controls(20.4,14.9)and(20.2,12.3)..(22.0,11.2)
    ..controls(24.0,9.8)and(28.1,10.0)..(28.85,9.0)
    ..controls(29.1,8.55)and(29.3,8.2)..(30.05,8.2);
  \draw[ft arrow,draw=ftGFPurple] (fOAC.west)..controls(20.4,14.9)and(20.2,12.3)..(22.0,11.2)
    ..controls(24.0,9.8)and(28.1,10.0)..(28.85,9.0)
    ..controls(29.1,8.55)and(29.3,8.2)..(30.05,8.2);
  \draw[ft underlay] (fOAD.south)..controls(23.0,11.6)and(26.0,9.8)..(28.85,9.0)
    ..controls(29.1,8.55)and(29.3,8.2)..(29.4,8.2);
  \draw[ft leader,draw=ftGFPurple] (fOAD.south)..controls(23.0,11.6)and(26.0,9.8)..(28.85,9.0)
    ..controls(29.1,8.55)and(29.3,8.2)..(29.4,8.2);
  \draw[ft underlay] (fOAB.east)..controls(27.5,12.77)and(28.5,12.7)..(28.6,11.1)
    ..controls(29.0,8.5)and(29.0,3.6)..(30.05,2.9);
  \draw[ft arrow,draw=ftGFRed] (fOAB.east)..controls(27.5,12.77)and(28.5,12.7)..(28.6,11.1)
    ..controls(29.0,8.5)and(29.0,3.6)..(30.05,2.9);

  % The new ray and the exact correspondence of the four old vertices.
  \path[rounded corners=4pt,fill=paperblue!40,draw=coreblue!22,line width=.8pt]
    (11.2,1.15) rectangle (27.85,9.1);
  \node[font=\fontsize{16}{18}\selectfont\bfseries,text=ink] at (19.525,8.38)
    {One new ray orbit, shared by the three tetrahedra};
  \node at (14.0,5.60) {\resizebox{3.65cm}{!}{% Generated by export_field_atlas_assets.jl from the maintained census.
\begingroup
\definecolor{g2TaxonBlue}{HTML}{1676A3}
\definecolor{g2TaxonYellow}{HTML}{E39A24}
\definecolor{g2HeartRed}{HTML}{D7443E}
\definecolor{g2EdgeBlue}{HTML}{205B7A}
\tikzset{
g2tree lettered vertex/.style={circle,fill=g2EdgeBlue,draw=white,line width=.7pt,inner sep=0pt,minimum size=5.0mm,font={\sffamily\fontsize{10}{10}\selectfont\bfseries},text=white},
  g2tree pendant/.style={draw=black!43,line width=.6pt,line cap=round},
  g2tree internal/.style={draw=g2EdgeBlue,line width=2pt,line cap=round},
  g2tree vertex/.style={circle,fill=g2EdgeBlue,draw=white,line width=.4pt,inner sep=0pt,minimum size=4.8pt},
  g2tree dark/.style={fill=g2TaxonBlue,draw=white,line width=.55pt,text=white,inner sep=0pt},
  g2tree light/.style={fill=g2TaxonYellow,draw=white,line width=.55pt,text=white,inner sep=0pt},
  g2tree weight/.style={circle,fill=white,draw=none,text=black,inner sep=.15pt,minimum size=10.0pt},
  g2tree aligned pendant/.style={fill=white,text=g2EdgeBlue,inner sep=.7pt,sloped,above=1pt,font={\sffamily\fontsize{10.25}{10.75}\selectfont}},
  g2tree aligned internal/.style={fill=white,text=g2EdgeBlue,inner sep=1pt,sloped,above=1pt,font={\sffamily\fontsize{10.25}{10.75}\selectfont\bfseries}},
}
\begin{tikzpicture}[x=1.0cm,y=1.0cm,baseline=(current bounding box.center)]
\coordinate (v1) at (1.65748,1.32180);
\coordinate (v2) at (0.91983,-1.91005);
\coordinate (v3) at (-2.06685,-0.47174);
\coordinate (v4) at (2.06685,-0.47174);
\coordinate (v5) at (-1.65748,1.32180);
\coordinate (v6) at (-0.91983,-1.91005);
\coordinate (v7) at (0.00000,2.12000);
\coordinate (v8) at (1.28094,0.38267);
\coordinate (v9) at (-0.07848,-1.10402);
\coordinate (v10) at (-1.23851,0.17965);
\coordinate (v11) at (0.32237,0.66755);
\coordinate (v12) at (-0.28632,-0.12585);
\draw[g2tree internal] (v8)--(v11);
\draw[g2tree internal] (v9)--(v12);
\draw[g2tree internal] (v10)--(v12);
\draw[g2tree internal] (v11)--(v12);
\draw[g2tree pendant] (v1)--(v8);
\draw[g2tree pendant] (v2)--(v9);
\draw[g2tree pendant] (v3)--(v10);
\draw[g2tree pendant] (v4)--(v8);
\draw[g2tree pendant] (v5)--(v10);
\draw[g2tree pendant] (v6)--(v9);
\draw[g2tree pendant] (v7)--(v11);
\node[g2tree weight,font={\sffamily\fontsize{10.25}{10.75}\selectfont}] at (0.80166,0.52511) {1};
\node[g2tree weight,font={\sffamily\fontsize{10.25}{10.75}\selectfont}] at (-0.18240,-0.61494) {1};
\node[g2tree weight,font={\sffamily\fontsize{10.25}{10.75}\selectfont}] at (-0.76241,0.02690) {1};
\node[g2tree weight,font={\sffamily\fontsize{10.25}{10.75}\selectfont}] at (0.01803,0.27085) {1};
\node[g2tree weight,font={\sffamily\fontsize{10.25}{10.75}\selectfont}] at (1.46921,0.85223) {0};
\node[g2tree weight,font={\sffamily\fontsize{10.25}{10.75}\selectfont}] at (0.42067,-1.50704) {0};
\node[g2tree weight,font={\sffamily\fontsize{10.25}{10.75}\selectfont}] at (-1.65268,-0.14604) {0};
\node[g2tree weight,font={\sffamily\fontsize{10.25}{10.75}\selectfont}] at (1.67389,-0.04454) {0};
\node[g2tree weight,font={\sffamily\fontsize{10.25}{10.75}\selectfont}] at (-1.44800,0.75073) {0};
\node[g2tree weight,font={\sffamily\fontsize{10.25}{10.75}\selectfont}] at (-0.49916,-1.50704) {0};
\node[g2tree weight,font={\sffamily\fontsize{10.25}{10.75}\selectfont}] at (0.16119,1.39378) {1};
\node[g2tree lettered vertex] at (v8) {A};
\node[g2tree lettered vertex] at (v9) {B};
\node[g2tree lettered vertex] at (v10) {C};
\node[g2tree lettered vertex] at (v11) {D};
\node[g2tree lettered vertex] at (v12) {E};
\node[g2tree dark,rectangle,minimum size=6.21mm,font={\sffamily\fontsize{12}{12}\selectfont\bfseries}] at (v1) {1};
\begin{scope}[shift={(v2)},scale=1.35,transform shape]
\path[fill=g2TaxonBlue,draw=white,line width=.55pt,line join=round] (0mm,3.0mm)--(-3.0mm,-2.2mm)--(3.0mm,-2.2mm)--cycle;
\node[font={\sffamily\fontsize{8.888888888888888}{8.888888888888888}\selectfont\bfseries},text=white,inner sep=0pt] at (0mm,-.35mm) {2};
\end{scope}
\node[g2tree dark,circle,minimum size=6.48mm,font={\sffamily\fontsize{12}{12}\selectfont\bfseries}] at (v3) {3};
\node[g2tree light,rectangle,minimum size=6.21mm,font={\sffamily\fontsize{12}{12}\selectfont\bfseries}] at (v4) {4};
\begin{scope}[shift={(v5)},scale=1.35,transform shape]
\path[fill=g2TaxonYellow,draw=white,line width=.55pt,line join=round] (0mm,3.0mm)--(-3.0mm,-2.2mm)--(3.0mm,-2.2mm)--cycle;
\node[font={\sffamily\fontsize{8.888888888888888}{8.888888888888888}\selectfont\bfseries},text=white,inner sep=0pt] at (0mm,-.35mm) {5};
\end{scope}
\node[g2tree light,circle,minimum size=6.48mm,font={\sffamily\fontsize{12}{12}\selectfont\bfseries}] at (v6) {6};
\begin{scope}[shift={(v7)},scale=1.35,transform shape]
\path[fill=g2HeartRed,draw=white,line width=.55pt]
  (0mm,-2.7mm)..controls(-.45mm,-2.1mm)and(-3.0mm,.35mm)..(-3.0mm,1.45mm)
  ..controls(-3.0mm,3.2mm)and(-.9mm,3.4mm)..(0mm,1.75mm)
  ..controls(.9mm,3.4mm)and(3.0mm,3.2mm)..(3.0mm,1.45mm)
  ..controls(3.0mm,.35mm)and(.45mm,-2.1mm)..(0mm,-2.7mm)--cycle;
\node[font={\sffamily\fontsize{8.888888888888888}{8.888888888888888}\selectfont\bfseries},text=white,inner sep=0pt] at (0mm,.35mm) {7};
\end{scope}
\end{tikzpicture}
\endgroup
}};
  \node[ft tree caption] at (14.0,3.38) {$O$: GF$(2)$ ray 3};
  \node[anchor=west,font=\fontsize{15}{17}\selectfont\bfseries,text=heartred] at (16.8,6.95)
    {GF$(2)$ rays $3,36,39$\quad (orbit size 3)};
  \node[anchor=west,align=left,text width=10.6cm,font=\fontsize{13}{16}\selectfont,text=ink]
    at (16.6,5.55) {In the quotient by lineality, the central ray\\is the sum of the four primitive vertex rays.};
  \node[anchor=west,align=left,text width=10.6cm,font=\fontsize{13}{16}\selectfont,text=muted]
    at (16.6,3.9) {The stabilizer fixes $A,B$ and exchanges $C,D$.\\The faces split $2+1+1$;\\the edges split $1+2+2+1$.};
  \node[font=\fontsize{12}{14}\selectfont,text=ink] at (19.525,2.32)
    {$\begin{array}{c|rrrr}\text{vertex}&A&B&C&D\\\hline
      \mathbb Q\text{ ray}&34&9&40&6\\
      \operatorname{GF}(2)\text{ ray}&1&2&13&14\end{array}$};
  \node[font=\fontsize{10.5}{12}\selectfont,text=muted] at (19.75,.28)
    {Representative Dressian tetrahedron 190. Orbit, cell and ray indices are field-specific census indices. The two tetrahedra are schematic projections of the same cell.};
\end{tikzpicture}
\endgroup
